% Fact check of the CBC Daybreak interview of May 20, 2026 with Stéphanie
% Valenzuela, borough mayor of Côte-des-Neiges–Notre-Dame-de-Grâce, on the
% Terrebonne bike path. Compile with pdflatex (twice), or run build.py.
\documentclass[11pt,letterpaper]{article}
\usepackage[T1]{fontenc}
\usepackage[utf8]{inputenc}
\usepackage[scaled=0.95]{helvet}
\renewcommand{\familydefault}{\sfdefault}
\usepackage[letterpaper,top=0.9in,bottom=0.95in,left=1in,right=1in]{geometry}
\usepackage{microtype}
\usepackage{enumitem}
\usepackage{titlesec}
\usepackage{fancyhdr}
\usepackage{xcolor}
\usepackage{booktabs}
\usepackage{longtable}
\usepackage{array}
\usepackage{textcomp}
\usepackage[hidelinks]{hyperref}
\urlstyle{same}
\newcommand{\docversion}{2026-10-26-15-27-55} % set by build.py (America/Toronto)
\hypersetup{pdfsubject={version \docversion}}

\setlength{\parindent}{0pt}
\setlength{\parskip}{0.5em}
\linespread{1.03}

\titleformat{\section}{\normalfont\large\bfseries}{\thesection}{0.6em}{}[\vspace{1pt}\titlerule]
\titlespacing*{\section}{0pt}{1.1em}{0.4em}
\titleformat{\subsection}{\normalfont\bfseries}{\thesubsection}{0.6em}{}
\titlespacing*{\subsection}{0pt}{0.8em}{0.25em}

\pagestyle{fancy}
\fancyhf{}
\renewcommand{\headrulewidth}{0pt}
\fancyfoot[C]{\footnotesize\color{gray}CBC Daybreak, May 20, 2026 -- fact check -- page \thepage}
\fancyfoot[R]{\tiny\color{gray}v.\,\docversion}

\newcommand{\repourl}{https://github.com/OpenMobilityData/Terrebonne-open-sources}
\newcommand{\repo}[1]{\href{\repourl}{#1}}
\newcommand{\doc}[2]{\href{\repourl/blob/main/#1}{#2}}
\newcommand{\arurl}{https://github.com/OpenMobilityData/Terrebonne-access-requests}
\newcommand{\pagesurl}{https://openmobilitydata.github.io/Terrebonne-access-requests}
\newcommand{\cliptime}[1]{\textsf{\small[#1]}}
\newcommand{\pgl}[2]{\href{\pagesurl/2026-581/raw/01_emails_mandate_counts.pdf\#page=#1}{\textsf{\small[2026-581, p.~#2]}}}
\newcommand{\pg}[1]{\pgl{#1}{#1}}
\newcommand{\exppg}[1]{\href{\repourl/blob/main/2023-exp-mobility-study/raw/EXP_Etude_de_mobilite_rue_de_Terrebonne_2023.pdf}{\textsf{\small[EXP 2023, p.~#1]}}}
\newcommand{\wc}[1]{\href{\repourl/blob/main/2020-2021-working-committee/raw/Comite_de_travail_recommandations_2021-06-13.pdf}{\textsf{\small[committee note, June 13, 2021, p.~#1]}}}
\newcommand{\quot}[1]{\begin{quote}\small #1\end{quote}}
\newcommand{\tr}[1]{\href{\repourl/blob/main/2026-10-19-city-council/derived/2026-10-19_03.md}{\textsf{\small[City council, Oct 19, 2026, #1]}}}
\newcommand{\octsix}[2]{\href{\repourl/blob/main/2026-10-06-borough-council/derived/#1}{\textsf{\small[Borough council, Oct 6, 2026, #2]}}}
\newcommand{\septeight}[1]{\href{\repourl/blob/main/2026-09-08-borough-council/}{\textsf{\small[Borough council, Sept 8, 2026, #1]}}}
\newcommand{\dbocttime}[1]{\href{\repourl/blob/main/2026-cbc-daybreak-interview/}{\textsf{\small[Daybreak, Oct 9, 2026, #1]}}}
\newcommand{\cjad}[1]{\href{\repourl/blob/main/2025-08-20-cjad-valenzuela/derived/2025-08-20_CJAD_Valenzuela_transcript_en.md}{\textsf{\small[CJAD, Aug 20, 2025, #1]}}}
\newcommand{\dbnov}[1]{\href{\repourl/blob/main/2025-11-06-cbc-daybreak-valenzuela/derived/2025-11-06_CBC_Daybreak_Valenzuela_transcript_en.md}{\textsf{\small[Daybreak, Nov 6, 2025, #1]}}}

\begin{document}

\begin{center}
{\Large\bfseries CBC \emph{Daybreak}, May 20, 2026: the borough mayor's statements on the Terrebonne bike path checked against the record}\\[0.4em]
{\large Fact check of the interview with Stéphanie Valenzuela}\\[0.6em]
Richard Hoge, Côte-des-Neiges--Notre-Dame-de-Grâce -- October 26, 2026\\
Published with the transcript at \repo{\repourl}
\end{center}

\section*{Summary}

On May~20, 2026, on her return from maternity leave, the borough mayor of
Côte-des-Neiges--Notre-Dame-de-Grâce, Stéphanie Valenzuela, was interviewed
on CBC Radio's \emph{Daybreak} (Montréal). Mobility is discussed from 6:43
and the Terrebonne bike path from 8:45 to 10:17. Those statements were
compared with the public record, in particular the file released under access
request 2026-581 (whose earliest document is dated three weeks after the
interview) and her own earlier interviews of August~20 and November~6, 2025.
In short:

\begin{enumerate}[nosep]
\item \textbf{« Something that is long term as a solution » \cliptime{9:00}.} Asked what residents should expect, the mayor said the borough was « really working on finding something that is long term as a solution ». This is the first public statement on the record that a lasting change to the 2024 layout was being sought. It is consistent with what the file later shows: on August~13, 2026 the mayor's chief of staff set a permanent mandate (a bidirectional lane on one side, parking restored on the other, \pg{1}), and on October~9 the councillor for Loyola said « for now, we're looking at it in a way that will make it permanent » \dbocttime{12:22}. What the public was told on October~6 was a seasonal measure and « a test » \octsix{2026-10-06_01.md}{40:49}.
\item \textbf{« Concrete scenarios that have been presented to us, and now we're just evaluating which ones we'll be putting in place » \cliptime{9:59}.} On May~20 scenarios already existed, before any outside study. The June~10, 2026 letter from the borough's engineering division to the City says the borough's engineers had already given a technical opinion on the two concepts sketched by the elected officials, whose objective was « de redonner la bande de stationnement qui avait été retranchée », and that the elected officials wanted « un avis officiel, neutre et externe » \pg{8}. On September~8 the councillor spoke of « a dozen or more of scenarios » \septeight{34:25} and on October~6 of « a dozen potential scenarios » \octsix{2026-10-06_05.md}{2:24:19}. No scenario has been published; the consultant's note and the City's analysis are withheld.
\item \textbf{« A lack of consultation » \cliptime{9:00}, « presented as a fait accompli » \cliptime{9:18}.} The same account as on CJAD on August~20, 2025 (« implemented it without consulting the residents », \cjad{3:34}). It leaves out the borough's working committee on mobility on Terrebonne (eight meetings, November~2020 to May~2021, with the street's institutions, school boards and residents), the EXP study published with the chosen scenario on October~30, 2023, the public information session of November~30, 2023, and the two council votes of February~2024, at the first of which she voted against. The host's remark that « that conversation had been going on since well before 2020 » was answered « Yeah, but ». Whether those steps amounted to consultation is a judgement; the check records what the statement leaves out.
\item \textbf{« We've taken the time over the last few months to really re-speak to all of the actors involved … but also to the bikers and to the pedestrians that do enjoy that bike path » \cliptime{9:41}.} The 508 pages released under request 2026-581 contain no record of any consultation, survey or meeting with residents or users. At the borough council of October~6, 2026, the Association des piétons et cyclistes de NDG asked why it had not been contacted \octsix{2026-10-06_02.md}{1:13:36}; the councillor answered that the association had been « en contact … avant, pendant, et juste avant que je vous parle » \octsix{2026-10-06_02.md}{1:21:34} and that there had been no « autres consultations publiques récemment » \octsix{2026-10-06_02.md}{1:15:02}; on October~9 he cited « 4,000 doors » and « more than 40 meetings » \dbocttime{01:53 and 02:37}. Which cyclists or pedestrians had been heard by May~20 is not on the record.
\item \textbf{« I think there is a middle ground » \cliptime{9:59}.} The frame recurs: on September~8 the councillor said « si on arrive à frustrer les deux côtés en même temps, on va savoir qu'on a trouvé une voie de passage au milieu » \septeight{2:38:50}, and on October~6 the mayor spoke of a « juste milieu » \octsix{2026-10-06_07.md}{50:53}. On the middle itself the record holds two technical findings: the City's street-design service found on June~17 that the borough's options « présentent plusieurs compromis importants sur le plan de la sécurité, du confort des usagers vulnérables et de la cohérence du réseau cyclable » \pg{6}; and the 2023 EXP study had already set aside a bidirectional lane, the long-term option of the August~13 mandate, as « pas souhaitable dans le contexte de la rue de Terrebonne » \exppg{51}.
\item \textbf{« We're actually coming up with a plan soon … bike paths … safety crossings » \cliptime{7:00}.} The borough has a mobility plan in force: the Plan local de déplacements 2024--2034, adopted unanimously on May~6, 2024 (CA24~170089) with her vote as councillor, with an internal review in 2027 and a public mid-term review in 2029 (\doc{2024-cdn-ndg-pld/}{documents}). On October~6 the councillor also announced « a longer term, more … sophisticated bike plan … for next year » \octsix{2026-10-06_01.md}{40:49}. Whether the plan announced is a new one or the 2027 review is not stated.
\item \textbf{The host's premise, « as part of the review of some bike paths, the city plans to look at one on Terrebonne » \cliptime{8:45}.} The City's release of February~10, 2026 names Terrebonne among the streets where the City works with the borough on safety, but its technical review covers four other lanes; on October~19, 2026 the executive committee member for mobility said those four reviews were not necessarily « exportable » to Terrebonne \tr{7:19:53}. The mayor did not take up the audit in her answer; it is the first of her three interviews in which the audit is not mentioned.
\end{enumerate}

\section{What was checked, and how}

\textbf{The interview.} CBC Radio One, \emph{Daybreak} (Montréal), clip
« Mayor Valenzuela is back in office from maternity leave », dated May~20,
2026, 11~min~41~s; host Sean Henry. The transcript published with this
document was drafted by machine and then checked passage by passage against
the audio. Only the statements on mobility and the Terrebonne bike path
(6:43 to 10:17) are examined; the Namur-Hippodrome site, the Cavendish
extension, parental leave and international outreach are outside this
check's scope.

\textbf{Sources.} Access request 2026-581 (bundle~01, cited by page); the
2020--2021 working committee documents; the EXP study of 2023 and its
release; the 2024 council records and the PLD 2024--2034; the City's release
of February~10, 2026; the borough councils of September~8 and October~6,
2026 and the City council of October~19, 2026 (transcripts checked by ear,
in this repository); the CJAD interview of August~20, 2025 and the
\emph{Daybreak} interviews of November~6, 2025 and October~9, 2026 (same).

\textbf{Terms.} \emph{Consistent} means the record supports the statement as
made; \emph{leaves out} that the record holds facts on the same point that
the statement does not mention; \emph{not on the record} that no document
bears on it.

\section{Findings in detail}

\subsection{« Long term as a solution »}

\cliptime{8:45} Host: « Still on the question of mobility, as part of the
review of some bike paths, the city plans to look at one on Terrebonne in
your borough. What should residents expect regarding specifically Terrebonne
and that bike path? » \cliptime{9:00} Mayor: « Yeah, what they should expect
is that we're really working on finding something that is long term as a
solution. What we inherited as a reality from the previous administration
was a lack of consultation, a lack of… »

\textbf{Consistent with the file as it later developed.} The earliest
document in the release under request 2026-581 is the June~10, 2026 letter
\pg{8}. The mandate from the mayor's chief of staff of August~13 is for a
permanent layout: « une piste cyclable bidirectionnelle d'un côté de la rue
de Terrebonne … qui respecte les normes minimales prévues par la loi » and
« le retour du stationnement d'un côté de la rue » \pg{1}; no winter measure
appears in it. The service request of August~20 adds « Volet~1 … mesure
transitoire pour l'hiver 2026-2027 » ahead of « Volet~2 », the permanent
lane \pg{3}. On October~6 the councillor announced the winter measure as
« seasonal » and « a test » \octsix{2026-10-06_01.md}{40:49}; on
October~9 he said « for now, we're looking at it in a way that will make it
permanent » \dbocttime{12:22}. The May~20 statement is therefore the
earliest public indication that the borough was seeking a lasting change,
five months before the winter measure was announced as a test.

\subsection{« A lack of consultation … a fait accompli »}

\cliptime{9:00} Host: « That had been going, that conversation had been
going on since well before 2020. » \cliptime{9:18} Mayor: « Yeah, but the
way that it was presented was there was no real openness to see what can be
modified, right? So I think it was really presented as a fait accompli. And
that was the grievance that was heightened really by the residents,
especially those that were directly impacted. If they lived on Terrebonne
and they were, you know, living the realization of this bike path, they just
didn't feel heard. »

\textbf{Leaves out the same steps as the August~2025 statement.} The
record before the 2024 lanes: the borough council's resolution CA20~170210
(2020) creating a working committee on mobility on Terrebonne, chaired by the
councillor for Notre-Dame-de-Grâce, with borough staff, the school boards,
the street's religious institutions, the HLM and residents, which met eight
times from November~19, 2020 to May~19, 2021 and recommended that the City
study a one-way scenario \wc{1}; the EXP mobility study commissioned by the
City and published by the borough on October~30, 2023 with the chosen
scenario, followed by exchanges with the institutional partners and a public
information session on November~30, 2023
(\doc{2023-exp-mobility-study/}{release}); the borough council vote of
February~5, 2024 (CA24~170018, 3--2, the speaker among the two against) and
the City council vote of February~19, 2024 (CM24~0179)
(\doc{2024-project-approval/}{documents}). On CJAD on August~20, 2025 she
had said the previous administration « implemented it without consulting the
residents » and that the working group's work was « completely …
disregarded » \cjad{3:34}; the committee's own note shows it recommended a
study, not a layout, and that the 2024 project made the street one-way as
the committee's unanimous recommendation had asked the City to study. That
residents on the street « didn't feel heard » is reported as their grievance
and is not checked here.

\subsection{« Re-speak to all of the actors involved »}

\cliptime{9:41} « And so we're taking the time and we've taken the time
over the last few months to really re-speak to all of the actors involved,
whether it's like the institutions, the church, the schools, the residents
that are lacking, you know, parking space, et cetera. But also to the bikers
and to the pedestrians that do enjoy that bike path. »

\textbf{Not on the record for cyclists and pedestrians.} The 508 pages
released on October~9, 2026 contain no minutes, survey, correspondence or
summary of any consultation with residents or users; the only exchanges
recorded are with the City's services, the fire department and the
consultant. At the borough council of October~6, 2026 a representative of
the Association des piétons et cyclistes de NDG asked « comment ça se fait
que vous n'aviez pas communiqué avec l'Association » before the announcement
\octsix{2026-10-06_02.md}{1:13:36}; the councillor answered that the
association had been « en contact … avant, pendant, et juste avant que je
vous parle » \octsix{2026-10-06_02.md}{1:21:34}, referring to a meeting held
in the days before the sitting, and said there had been no « autres
consultations publiques récemment », the file resting on « énormément de
dossiers, de données, de rapports » \octsix{2026-10-06_02.md}{1:15:02}. On
October~9 he cited « 4,000 doors » knocked on during the campaign and
« more than 40 meetings » \dbocttime{01:53 and 02:37}, and on October~6
« more than 30 meetings on this topic with the services »
\octsix{2026-10-06_01.md}{38:16}. Which cyclists or pedestrians had been
heard by May~20, 2026, and what they said, appears in no document.

\subsection{« A middle ground »; « concrete scenarios »}

\cliptime{9:59} Host: « Is there a middle ground in that conversation,
though? » Mayor: « I think there is a middle ground, and that's what we're
working on with our services and our engineers. We have concrete scenarios
that have been presented to us, and now we're just evaluating which ones
we'll be putting in place. »

\textbf{Scenarios before the study.} Three weeks after the interview, the
June~10 letter states that the elected officials wished to revisit the 2024
layout « avec l'objectif de redonner la bande de stationnement qui avait été
retranchée », that the lanes' width was « jugée disproportionnée par rapport
au volume de cyclistes observé … selon les élus », that the borough's
engineers had already issued a technical opinion on the two concepts the
elected officials had sketched, and that the elected officials wanted « un
avis officiel, neutre et externe » \pg{8}. The City's street-design service
answered on June~17; the borough's summary of June~22 says the options
« présentent plusieurs compromis importants sur le plan de la sécurité, du
confort des usagers vulnérables et de la cohérence du réseau cyclable »
\pg{6}; the analysis is withheld. The outside study was ordered on
August~20 \pg{3}. On September~8 the councillor said the borough was « at
the final stage of the study » with « a dozen or more of scenarios »
\septeight{34:25}; on October~6, « a dozen potential scenarios »
\octsix{2026-10-06_05.md}{2:24:19}. The scenarios, the consultant's note
and the City's analysis are withheld under section~39 of the Access Act,
« until a decision is made on the recommendation ».

\textbf{The middle, as framed later.} On September~8: « si on arrive à
frustrer les deux côtés en même temps, on va savoir qu'on a trouvé une voie
de passage au milieu » \septeight{2:38:50}. On October~6 the mayor: « des
gens qui veulent qu'on l'enlève carrément … ce n'est pas où nous voulons
aller », a « juste milieu » \octsix{2026-10-06_07.md}{50:53}. On the
technical content of the middle, the record holds the June~17 finding above
and, for the long-term option of the August~13 mandate, the 2023 EXP
study's position that « un aménagement bidirectionnel n'est pas souhaitable
dans le contexte de la rue de Terrebonne … notamment avec la présence de
nombreuses entrées charretières et avec la gestion de plusieurs
intersections » \exppg{51}, none of its three concepts being a two-way lane.

\subsection{« Coming up with a plan soon »}

\cliptime{6:43} Host: « What is the borough's general plan when it comes to
mobility and stuff like that in Côte-des-Neiges-NDG? » \cliptime{7:00}
Mayor: « So when it comes to mobility and stuff like that in NDG, we're
actually coming up with a plan soon that we're going to be working on to
really integrate connectivity, whether it's for bike paths or for, you know,
safety crossings that are important in our borough. … And we hope that we'll
be able to announce something very soon. »

\textbf{A plan is in force.} The borough's Plan local de déplacements
2024--2034 was adopted unanimously on May~6, 2024 (CA24~170089, moved by the
then mayor, with the speaker present and voting for); it provides an
internal review in 2027 and a public mid-term review in 2029, and measure
3.1.1 upgrades « chaussées désignées » to protected facilities
(\doc{2024-cdn-ndg-pld/}{documents}). It does not name Terrebonne. On
October~6 the councillor announced « a longer term, more … sophisticated
bike plan … for next year » \octsix{2026-10-06_01.md}{40:49}. Whether the
plan announced on May~20 is the 2027 review or a new document is not stated;
nothing had been published by October~26, 2026.

\subsection{The three interviews side by side}

\begin{small}
\begin{longtable}{>{\raggedright\arraybackslash}p{2.6cm}>{\raggedright\arraybackslash}p{6.2cm}>{\raggedright\arraybackslash}p{6.0cm}}
\toprule
Date & Said & Record \\
\midrule
\endhead
Aug~20, 2025 (candidate) & an audit of the path's « security … implementation … impact on local residents » \cjad{3:34}; « without consulting the residents » & — \\
Nov~6, 2025 (mayor-elect) & « we wait for the results. We keep what is there as a network » \dbnov{7:55}; an audit of our own if the existing one is not independent \dbnov{8:54} & Nov~10: City note, no safety audit of the network (2025-12847) \\
Feb~10, 2026 & — & City review of four lanes launched; results at the end of 2026; Terrebonne named \\
May~20, 2026 (mayor) & « something that is long term as a solution » \cliptime{9:00}; « a lack of consultation » \cliptime{9:00}; « concrete scenarios … presented to us » \cliptime{9:59}; « a middle ground »; the audit not mentioned & — \\
Jun~10 -- 22, 2026 & — & objective: parking restored; engineers' opinion already given on the elected officials' two concepts; City analysis: « plusieurs compromis importants … sécurité », withheld (2026-581) \\
Aug~13 -- 20, 2026 & — & permanent mandate (bidirectional, parking); winter removal added as « mesure transitoire » \\
Sept~8, 2026 & councillor: « a dozen or more of scenarios »; « frustrer les deux côtés » & — \\
Oct~6, 2026 & orientation: one direction to parking each winter, « a test »; mayor: « juste milieu » & before the City's review reported \\
Oct~9, 2026 & councillor: « for now … permanent » & — \\
Oct~19, 2026 & — & City (DeSousa): results « may not be before » the vote; the four reviews not « exportable » \\
\bottomrule
\end{longtable}
\end{small}

\subsection{Other statements}

\cliptime{7:16} « We're one borough that is surrounded by other
municipalities, whether it's Côte Saint-Luc, Hampstead, Montreal West »:
geography, consistent (Westmount is named at 8:19).

\cliptime{3:15} « my colleague, Sonny Moroz, who was interim mayor » and
\cliptime{2:03} « At the municipal level, we get 18 weeks »: not examined.

\cliptime{6:43} Host: the decision « to scale back plans to add parking
meters on and around Sherbrooke Street West »: outside this check's scope.

\section{What this check does not establish}

It compares statements with records. It does not establish why any statement
was made and draws no conclusion about intent. The scenarios, the
consultant's note and the City's analysis of June~2026 are withheld; the
City's review had not reported when this check was written.

\section{Sources}

\begin{small}
\begin{longtable}{>{\raggedright\arraybackslash}p{4.2cm}>{\raggedright\arraybackslash}p{10.6cm}}
\toprule
Source & Where \\
\midrule
\endhead
Interview & CBC \emph{Daybreak}, May~20, 2026: \url{https://www.cbc.ca/listen/live-radio/1-15-daybreak-montreal/clip/16215928}; transcript: \doc{2026-05-20-cbc-daybreak-valenzuela/derived/2026-05-20_CBC_Daybreak_Valenzuela_transcript_en.md}{this folder} \\
Earlier interviews & \doc{2025-08-20-cjad-valenzuela/}{CJAD, August~20, 2025}; \doc{2025-11-06-cbc-daybreak-valenzuela/}{Daybreak, November~6, 2025}; \doc{2026-cbc-daybreak-interview/}{Daybreak, October~9, 2026 (the councillor)} \\
Request 2026-581 & Borough response of October~9, 2026, bundle 01: \href{\pagesurl/2026-581/raw/01_emails_mandate_counts.pdf}{PDF} (pp.~1--4, 6, 8); summaries in the \href{\arurl/tree/main/2026-581}{access-requests repository} \\
Working committee & \doc{2020-2021-working-committee/}{CA20~170210; note of June~13, 2021; file 1215284005} \\
EXP study & \doc{2023-exp-mobility-study/}{Study (2023) and the release of October~30, 2023} \\
2024 approvals; PLD & \doc{2024-project-approval/}{CA24~170018, CM24~0179, CA24~170131}; \doc{2024-cdn-ndg-pld/}{PLD 2024--2034, CA24~170089} \\
City review & \doc{2026-city-audit-and-plan-velo/}{Release of February~10, 2026 and the Réalisons Montréal timeline} \\
Borough councils & \doc{2026-09-08-borough-council/}{September~8, 2026} and \doc{2026-10-06-borough-council/}{October~6, 2026}, transcripts checked by ear \\
City council, Oct~19, 2026 & \doc{2026-10-19-city-council/}{Transcripts checked by ear} (clip 03) \\
\bottomrule
\end{longtable}
\end{small}

\end{document}
